% from ega2-1-fr.tex
\oldpage[II]{192}
quasi-cohérent associé à $\mathcal{M}$
(\hyperref[II.1.4.3-fr]{1.4.3}) lorsque $\mathcal{M}$ est considéré comme
$\mathcal{S}$-Module \emph{non gradué}, et par
$\operatorname{Proj}_0(\mathcal{M})$ le $\mathcal{O}_X$-Module
quasi-cohérent associé à $\mathcal{M}$, $\mathcal{M}$ étant cette fois
considéré comme $\mathcal{S}$-Module gradué (autrement dit le
$\mathcal{O}_X$-Module noté $\widetilde{\mathcal{M}}$ dans
(\hyperref[II.3.2.2-fr]{3.2.2})). Nous poserons de plus
\begin{equation}
\label{II.8.12.1.1-fr}
  \mathcal{M}_X=\operatorname{Proj}(\mathcal{M})
  =\bigoplus_{n\in\mathbf{Z}}\operatorname{Proj}_0(\mathcal{M}(n))~;
\tag{8.12.1.1}
\end{equation}
la $\mathcal{O}_X$-Algèbre graduée quasi-cohérente $\mathcal{S}_X$ étant
définie par (\hyperref[II.8.6.1.1-fr]{8.6.1.1}),
$\operatorname{Proj}(\mathcal{M})$ est muni d'une structure de
$\mathcal{S}_X$-Module gradué (quasi-cohérent), au moyen des homomorphismes
canoniques (\hyperref[II.3.2.6.1-fr]{3.2.6.1})
\begin{equation}
\label{II.8.12.1.2-fr}
  \mathcal{O}_X(m)\otimes_{\mathcal{O}_X}
  \operatorname{Proj}_0(\mathcal{M}(n))
  \to \operatorname{Proj}_0(\mathcal{S}(m)\otimes_{\mathcal{S}}\mathcal{M}(n))
  \to \operatorname{Proj}_0(\mathcal{M}(m+n))
\tag{8.12.1.2}
\end{equation}
la vérification des axiomes des Modules se faisant à l'aide du diagramme
commutatif (\hyperref[II.2.5.11.4-fr]{2.5.11.4}).

Si $Y=\operatorname{Spec}(A)$ est affine, $\mathcal{S}=\widetilde{S}$ et
$\mathcal{M}=\widetilde{M}$, où $S$ est une $A$-algèbre graduée et $M$ un
$S$-module gradué, alors, pour tout élément homogène $f\in S_+$, on a
\begin{equation}
\label{II.8.12.1.3-fr}
  \Gamma(X_f,\operatorname{Proj}(\widetilde{M}))=M_f
\tag{8.12.1.3}
\end{equation}
en vertu des définitions et de (\hyperref[II.8.2.9.1-fr]{8.2.9.1}).

Considérons maintenant le $\widehat{\mathcal{S}}$-Module gradué
quasi-cohérent
\begin{equation}
\label{II.8.12.1.4-fr}
  \widehat{\mathcal{M}}
  =\mathcal{M}\otimes_{\mathcal{S}}\widehat{\mathcal{S}}
\tag{8.12.1.4}
\end{equation}
($\widehat{\mathcal{S}}$ étant défini par
(\hyperref[II.8.3.1.1-fr]{8.3.1.1}))~; on en déduit un
$\mathcal{O}_{\widehat{C}}$-Module gradué quasi-cohérent
$\operatorname{Proj}_0(\widehat{\mathcal{M}})$, que nous noterons aussi
\begin{equation}
\label{II.8.12.1.5-fr}
  \mathcal{M}^{\square}=\operatorname{Proj}_0(\widehat{\mathcal{M}}).
\tag{8.12.1.5}
\end{equation}
Il est clair (\hyperref[II.3.2.4-fr]{3.2.4}) que
$\mathcal{M}^{\square}$ est un foncteur additif exact en $\mathcal{M}$,
commutant aux sommes directes et aux limites inductives.
\end{env}

\begin{proposition}[8.12.2]
\label{II.8.12.2-fr}
Avec les notations de (\hyperref[II.8.3.2-fr]{8.3.2}), on a des
isomorphismes canoniques fonctoriels
\begin{equation}
\label{II.8.12.2.1-fr}
  i^*(\mathcal{M}^{\square})\xrightarrow{\sim}\widetilde{\mathcal{M}},
  \qquad
  j^*(\mathcal{M}^{\square})\xrightarrow{\sim}
  \operatorname{Proj}_0(\mathcal{M}).
\tag{8.12.2.1}
\end{equation}
\end{proposition}

\begin{proof}
En effet, $i^*(\mathcal{M}^{\square})$ s'identifie canoniquement à
$\left(\widehat{\mathcal{M}}/(z-1)\widehat{\mathcal{M}}\right)^{\sim}$
sur
$\operatorname{Spec}(\widehat{\mathcal{S}}/(z-1)\widehat{\mathcal{S}})$
en vertu de (\hyperref[II.3.2.3-fr]{3.2.3})~; le premier des isomorphismes
canoniques (\hyperref[II.8.12.2.1-fr]{8.12.2.1}) se déduit alors aussitôt
(\hyperref[II.1.4.1-fr]{1.4.1}) de l'isomorphisme canonique
$\widehat{\mathcal{M}}/(z-1)\widehat{\mathcal{M}}
\xrightarrow{\sim}\mathcal{M}$. D'autre part, l'immersion canonique
$j:X\to\widehat{C}$ correspond à l'homomorphisme canonique
$\widehat{\mathcal{S}}\to\mathcal{S}$ de noyau
$z\widehat{\mathcal{S}}$ (\hyperref[II.8.3.2-fr]{8.3.2})~; le second
homomorphisme (\hyperref[II.8.12.2.1-fr]{8.12.2.1}) est le cas particulier
de l'homomorphisme canonique (\hyperref[II.3.5.2-fr]{3.5.2, (ii)}), du fait
que l'on a ici
$\widehat{\mathcal{M}}\otimes_{\widehat{\mathcal{S}}}\mathcal{S}
=\mathcal{M}$~; pour vérifier que c'est un isomorphisme, on peut se borner au
cas où $Y=\operatorname{Spec}(A)$ est affine,
$\mathcal{S}=\widetilde{S}$, $\mathcal{M}=\widetilde{M}$~; en se reportant à
(\hyperref[II.2.8.8-fr]{2.8.8}), la vérification que pour tout $f$ homogène
dans $S_+$, l'homomorphisme précédent, restreint à $X_f$, se réduit à un
isomorphisme, est alors immédiate.
\end{proof}

